\documentclass[11pt]{article}
\usepackage{amsmath,amssymb,amsthm}
\usepackage[margin=1in]{geometry}
\newtheorem{proposition}{Proposition}
\title{Three intersection applications for a fixed Fi22 selector}
\author{Mathematical intake working note}
\date{}
\begin{document}
\maketitle

This note applies the existing theorem on the intersection property in rank
three to three already accepted generation records. It introduces no group
calculation or general theorem. The named group, original membership, action,
class, Reader, two-fibre and $H=B$ interpretations remain the qualified
premises of those generation records. In particular, $G$ is the named inner
simple group $\mathrm{Fi}_{22}$, not its full automorphism group.

\section*{The three records}

The exact records are specified by their free indices, owners and full
witness-row digests in \texttt{proposal.json}. Their short identifiers are
\[
\begin{array}{c|r|l}
\text{free index}&\text{owner}&\text{witness digest prefix}\\\hline
265&474&\texttt{884cc09ae6be3a98}\\
270&512&\texttt{3446a5734cc79fe8}\\
271&514&\texttt{5d171be3dfa04919}
\end{array}
\]
The canonical batch with digest beginning \texttt{e44daa80} registers all
three generation facts. Its strict result, \texttt{dc5998e7}, and distinct
actual review, \texttt{82561a7c}, list exactly these identities and assert
\begin{equation}\label{generation}
                 \langle ab_i,b_ic\rangle=G.
\end{equation}
Their generation proof uses already checked supplied cycles and an
attributed maximal-subgroup argument. Neither calculation is repeated here.
The durable retained copies are joined to the original descriptors by the
accepted retention manifest \texttt{50336071}; the inputs record those
complete aliases rather than depending on a surviving temporary pathname.

The recorded current status of each of these rows is $2$: generation has
been accepted, while its separate intersection property has not been
registered. An old zero intersection-credit field in the strict result is
left unchanged. The present note supplies a separate mathematical deduction.

\section*{Hypotheses and deduction}

The existing note \emph{The intersection property in rank three for simple
groups}, SHA256
\texttt{0b940e5a217a5e54e91fd13617fd822dd4daebfac6cbdacbd17e0950070c9216},
proves the following statement: if a finite nonabelian simple group $G$ is
generated by nonidentity involutions $a,b,c$, where $a,c$ are distinct and
commute, then the triple has the full intersection property. Its independent
review, \texttt{b7610c40}, records the precise original-domain hypotheses of
the previous application. There is no hypothesis about adjacent product
orders or coprimality. We reuse the elementary proof and its existing
attribution; no historical novelty is claimed.

\begin{proposition}
Under the same qualified original-domain premises as their accepted
generation records, the three triples $(a,b_i,c)$, for
$i\in\{265,270,271\}$, have the full intersection property in rank three.
\end{proposition}
\begin{proof}
The whole-operand decoder with digest beginning \texttt{60eedb3c} checks
that $a,c$ are nonidentity involutions, that $a\ne c$, and that $ac=ca$.
The distinct actual review of this strict invocation closes that Source
and its inputs. These endpoint checks already occurred and are not repeated
by this note.

The accepted membership interpretation places the seed involution $d_0$
in $G$ and imports $G\lhd A$. Every specified ambient conjugate of the
seed therefore belongs to $G$, without asserting that $A\le G$. The
accepted original endpoint identifications put $a,c$ in $G$. The exact
original class and two-fibre interpretation identifies the product
representative attached to each bound witness with its intended middle
involution $b_i$. Thus $b_i\in G$ and $b_i^2=1$ are precisely the borrowed
class and membership premises already used in review \texttt{b7610c40}.
They are not derived here from an assertion that the three internal partner
factors commute. No such assertion is needed.

It follows from \eqref{generation} and these membership premises that
\[
 G=\langle ab_i,b_ic\rangle
       \le\langle a,b_i,c\rangle\le G.
\]
The three involutions therefore generate $G$. Moreover $b_i\ne1$; if
$b_i=1$, equation~\eqref{generation} would give
$G=\langle a,c\rangle$, an abelian group, contrary to the named nonabelian
group $G$. All hypotheses of the existing theorem are satisfied. Writing
$r_0=a$, $r_1=b_i$, $r_2=c$, it gives
\[
 \langle r_j:j\in I\rangle\cap\langle r_j:j\in J\rangle
       =\langle r_j:j\in I\cap J\rangle
       \qquad(I,J\subseteq\{0,1,2\}).
\]
Finally, the accepted witness, class and two-fibre interpretations attach
this triple to the same original owner as the generation record. The
inherited $H=B$ scope is retained. Conjugation by the accepted $H$ induces
automorphisms of $G$ and preserves subgroup intersections. It therefore
preserves the intersection property, even if $H$ contains outer
automorphisms. The integer owner label is not interpreted as an element or
a word of $H$.
\end{proof}

\section*{Scope and attribution}

The common theorem, its proof and review, the existing generation
certificates and controls, and the original model and class meanings are
borrowed. The contribution here is the explicit application to these three
records with their exact identities and hypothesis joins. No new order,
matrix product, point image, witness evaluation or search is used. This
note neither re-establishes the imported model recognition nor enlarges the
accepted original domain.

This is a proposal awaiting independent application review and separate
Root issuance. It makes no canonical change. If subsequently adopted, the
three status bytes would change from $2$ to $6$. The separate intersection
count would rise from $721$ to $724$; generation would remain $724$,
nongeneration $744$, settled rows $1468$, and unpaid rows $27440$. Those
prospective counts are not current counts. No original manuscript,
conjecture, Lean proof or compression measurement receives new credit.
\end{document}
